\input{preamble}

% ---- GROUP MACROS (begin): your own \newcommand lines go here, and only here ----
\newcommand{\Mean}{\bar{x}}
% ---- GROUP MACROS (end) ----

\title{A sample mean in one pass}
\author{First Member \and Second Member}
\date{ScaDaMaLe, WASP 2026 \quad (an own-paper fixture)}

\begin{document}
\maketitle

\begin{abstract}
A running mean that keeps one number of state, and the proof that it ends at the sample mean.
\end{abstract}

\section{Setting}\label{sec-setting}

\begin{definition}\label{Dfn-SampleMean}
For $x_1,\dots,x_n \in \mathbb{R}$ with $n \ge 1$, the \emph{sample mean} is
$\Mean = \frac{1}{n} \sum_{i=1}^{n} x_i$.
\end{definition}

\begin{lemma}\label{Lem-OnePass}
Let $m_0 = 0$ and $m_k = m_{k-1} + (x_k - m_{k-1})/k$ for $k = 1,\dots,n$. Then $m_n = \Mean$.
\end{lemma}

\begin{proof}
By induction on $k$, $m_k = \frac{1}{k} \sum_{i=1}^{k} x_i$: it holds for $k = 1$, since $m_1 = x_1$, and if it holds
for $k - 1$, then $m_k = \bigl((k-1)\, m_{k-1} + x_k\bigr)/k = \frac{1}{k} \sum_{i=1}^{k} x_i$. Take $k = n$.
\end{proof}

\section{Discussion}\label{sec-discussion}

\Cref{Lem-OnePass} computes the mean of \cref{Dfn-SampleMean} in one pass over the data, with no running sum that
can grow without bound (\cref{Dfn-SampleMean,Lem-OnePass}). In an exponential family the mean parameter is the
gradient of the log-partition function~\cite[Prop.~3.1]{wainwright2008graphical}, so this pass is also the moment
step of a maximum-likelihood fit (\cref{sec-setting}).

\section*{Authors' Contributions}

Each member writes their own line.

\section*{Use of AI tools}

None.

\bibliographystyle{alpha}
\bibliography{refs}

\end{document}
